\chapter{Statistics Hardening: Confidence Intervals, Significance Tests, and a Minimum-n Floor (Phase 3.1)}
\label{chap:p3-1}

\section{Goal}

Every headline number the pilot reported was a bare point estimate: 36.2\% descriptive accuracy,
28.3\% hidden drift, and so on, with no interval and no significance test behind the comparisons. The
plan makes this a gating requirement --- attach bootstrap confidence intervals and Wilcoxon signed-rank
tests (matching Prof.\ Abdallah's prior IDS work) to every headline metric, and enforce a minimum
samples-per-class floor before a rate is quoted as a stable estimate --- and states it must be in place
before any number is carried into the paper. This phase builds those three instruments and applies them
to the recorded results.

\section{What Was Done}

A pure, unit-tested \texttt{stats} module provides three functions: \texttt{bootstrap\_ci} (percentile
bootstrap, NaN-dropping, reproducible for a fixed seed, zero-width on a constant sample);
\texttt{wilcoxon\_signed\_rank} (paired nonparametric test, returning $p=1.0$ rather than raising when
every difference is zero); and \texttt{meets\_min\_n} against a floor of 30. A report script
(\texttt{13\_statistics\_report.py}) applies them to the headline metrics from the existing result CSVs
--- no new model calls --- and writes \texttt{statistics\_summary.csv}. Nine unit tests cover the
bootstrap (point, bracketing, reproducibility, constant, NaN, empty), the Wilcoxon (signal, all-ties),
and the floor.

\section{Results}

\subsection{Headline metrics now carry intervals, and the weak ones are flagged}

\begin{center}
\begin{tabular}{p{5.0cm}p{4.6cm}p{2.6cm}}
\toprule
\textbf{metric} & \textbf{estimate [95\% CI]} & \textbf{n} \\
\midrule
descriptive accuracy (top-1)          & 36.2\% [28.8, 43.6] & 163 \\
\ \ compliant                          & 28.0\% [16.0, 40.0] & 50 \\
\ \ fall\_hazard                       & 46.0\% [32.0, 60.0] & 50 \\
\ \ ppe\_violation                     & 28.0\% [16.0, 40.0] & 50 \\
\ \ struck\_by\_risk                   & 61.5\% [38.5, 84.6] & \textbf{13 (below floor)} \\
multi-hazard completeness (complete)  & 27.3\% [9.1, 45.5]  & \textbf{22 (below floor)} \\
\bottomrule
\end{tabular}
\captionof{table}{Bootstrap 95\% confidence intervals for headline rates. The two subgroups below the
minimum-n floor of 30 --- struck\_by\_risk ($n{=}13$) and the multi-hazard set ($n{=}22$) --- carry
intervals so wide ($\pm 23$ and $\pm 18$ points) that their point estimates cannot be quoted as stable,
and the report says so.}
\label{tab:p31-ci}
\end{center}

The intervals do real work immediately. The struck\_by\_risk descriptive accuracy (61.5\%) looked like
the strongest class in the pilot; its interval, $[38.5, 84.6]$ on $n{=}13$, shows that apparent strength
is indistinguishable from noise. The multi-hazard completeness rate from Phase~2.5 (27.3\%) is confirmed
as an existence result rather than a precise rate: $[9.1, 45.5]$. Both are flagged below the floor,
which is precisely the case the multi-label recovery (Phase~1.2) and the scale run are meant to fix ---
the statistics make the under-powering quantitative rather than a hand-wave.

\subsection{Two Phase-2 findings are statistically significant}

Paired Wilcoxon tests confirm that two of the phase-2 corrections reflect real signal, not sampling
noise:

\begin{center}
\begin{tabular}{p{6.4cm}p{3.0cm}p{2.6cm}}
\toprule
\textbf{paired comparison} & \textbf{effect} & \textbf{$p$} \\
\midrule
targeted vs.\ random-location occlusion drift & 0.198 vs.\ 0.142 & $1.1\times10^{-3}$ \\
stability agreement, temp 0.3 vs.\ 1.0        & 0.871 vs.\ 0.797 & $1.8\times10^{-2}$ \\
\bottomrule
\end{tabular}
\captionof{table}{Wilcoxon signed-rank tests on paired per-image measurements. The targeted-vs-random
occlusion distinction (Phase~2.4) and the temperature dose-response (Phase~2.3) are both significant at
$\alpha = 0.05$.}
\label{tab:p31-wilcoxon}
\end{center}

That targeted occlusion drifts the explanation significantly more than a random-location occlusion of
matched area ($p = 1.1\times10^{-3}$) puts a significance test behind Phase~2.4's central claim, and the
significant temperature effect ($p = 1.8\times10^{-2}$) does the same for Phase~2.3's dose-response.
These are the first significance tests in the project; the pilot compared such numbers by eye.

\section{Standard vs. Non-Standard Choices}

\textbf{Standard.} Bootstrap intervals, nonparametric paired tests, and a minimum-n floor are baseline
empirical-ML statistics; their absence in the pilot was the gap.

\textbf{Non-standard, and the contribution.} The floor is wired to \emph{flag} rather than silently
drop: struck\_by\_risk and the multi-hazard set still appear, with their wide intervals and an explicit
``below floor'' mark, so the reader sees both the estimate and its unreliability rather than a
suspiciously clean table with the hard cases quietly removed. And the significance tests are applied to
the project's \emph{own} phase-2 findings, holding them to the same bar the pilot's numbers will be held
to --- confirming two of them rather than assuming them.

\section{Implications}

Every headline number can now be reported as an estimate with an interval, every paired comparison with
a $p$-value, and every under-powered subgroup with an explicit flag. The two below-floor cases
(struck\_by\_risk, multi-hazard) are the concrete statistical targets the scale run must hit by drawing
more samples --- the same $n$-growth Phases~1.2 and~2.5 set up. The \texttt{stats} helpers are the
building blocks the programmatic report generation (Phase~3.4) will call for every table, so no number
reaches the paper without its interval.
